\chapter{Studio della Complessità}

In questo capitolo viene analizzata la complessità temporale e spaziale 
delle operazioni principali della struttura \texttt{HashRBTree}.

\section{Complessità Temporale}

\subsection{Funzione di Hash}

La funzione \texttt{hashFunction} itera su tutti i caratteri della chiave 
eseguendo un numero costante di operazioni per ciascuno. La complessità è:

\begin{equation}
    T_{hash} = O(m)
\end{equation}

dove $m$ è la lunghezza della chiave in caratteri.

\subsection{Ricerca nell'Albero}

Il metodo \texttt{searchNode} naviga l'albero Red-Black dall'alto verso il 
basso confrontando il valore \texttt{id} ad ogni nodo. Poiché l'albero è 
bilanciato per costruzione, l'altezza è garantita essere $O(\log n)$, dove 
$n$ è il numero di nodi distinti nell'albero. La complessità è:

\begin{equation}
    T_{searchNode} = O(\log n)
\end{equation}

\subsection{Ricerca nella Lista di Collisione}

Una volta individuato il nodo tramite \texttt{searchNode}, il metodo 
\texttt{search} scorre la lista \texttt{items} per trovare la chiave esatta. 
Indicando con $k$ il numero di coppie in collisione su quel nodo:

\begin{equation}
    T_{search} = O(m + \log n + k)
\end{equation}

Nel caso medio, con una buona funzione di hash, $k$ è costante e la 
complessità si riduce a $O(m + \log n)$.

\subsection{Inserimento}

L'operazione \texttt{insert} esegue in sequenza:

\begin{enumerate}
    \item Calcolo dell'hash: $O(m)$
    \item Ricerca del nodo: $O(\log n)$
    \item In caso di collisione, aggiunta in coda alla lista: $O(k)$
    \item In caso di nuovo nodo, inserimento BST: $O(\log n)$
    \item Ripristino proprietà RBT con \texttt{insertFix}: $O(\log n)$
\end{enumerate}

La complessità complessiva è:

\begin{equation}
    T_{insert} = O(m + \log n)
\end{equation}

\texttt{insertFix} esegue al più $O(\log n)$ recoloring risalendo l'albero, 
e al più 2 rotazioni che sono operazioni $O(1)$.

\subsection{Cancellazione}

L'operazione \texttt{remove} esegue in sequenza:

\begin{enumerate}
    \item Calcolo dell'hash: $O(m)$
    \item Ricerca del nodo: $O(\log n)$
    \item Rimozione dalla lista: $O(k)$
    \item Se la lista è vuota, rimozione dall'albero: $O(\log n)$
    \item Ripristino proprietà RBT con \texttt{deleteFix}: $O(\log n)$
\end{enumerate}

La complessità complessiva è:

\begin{equation}
    T_{remove} = O(m + \log n)
\end{equation}

\texttt{deleteFix} esegue al più $O(\log n)$ recoloring e al più 3 rotazioni.

\subsection{Riepilogo}

\begin{center}
\begin{tabular}{|l|c|c|}
\hline
\textbf{Operazione} & \textbf{Caso Medio} & \textbf{Caso Peggiore} \\
\hline
\texttt{hashFunction} & $O(m)$ & $O(m)$ \\
\hline
\texttt{search}       & $O(m + \log n)$ & $O(m + \log n + k)$ \\
\hline
\texttt{insert}       & $O(m + \log n)$ & $O(m + \log n + k)$ \\
\hline
\texttt{remove}       & $O(m + \log n)$ & $O(m + \log n + k)$ \\
\hline
\end{tabular}
\end{center}

\vspace{0.5cm}
dove $n$ è il numero di nodi nell'albero, $m$ è la lunghezza della chiave 
e $k$ è il numero di collisioni sul nodo.

\section{Complessità Spaziale}

La struttura occupa in memoria:

\begin{itemize}
    \item $O(n)$ per i nodi dell'albero, dove $n$ è il numero di chiavi 
    distinte per hash;
    \item $O(n + p)$ considerando anche le coppie nelle liste di collisione, 
    dove $p$ è il numero totale di coppie inserite;
    \item $O(\log n)$ per lo stack delle chiamate ricorsive di 
    \texttt{printHelper} e \texttt{destroyTree}.
\end{itemize}

La complessità spaziale complessiva è:

\begin{equation}
    S = O(n + p)
\end{equation}